% =====================================================================
\chapter{Aplicações das derivadas}

Com as regras de derivação em mãos, passamos às suas consequências. A derivada responde a
perguntas muito diferentes: com que rapidez uma grandeza varia quando outra varia (taxas
relacionadas), qual é o maior ou o menor valor de uma função (otimização), qual é a forma de
um gráfico, como calcular limites difíceis (regra de L'Hôpital) e como aproximar raízes de
equações (método de Newton). Terminamos com o caminho inverso, as antiderivadas, que abrem
as portas para o capítulo de integração.

% ---------------------------------------------------------------------
\section{Taxas relacionadas}

Quando duas grandezas estão ligadas por uma equação e ambas variam com o tempo, suas taxas
de variação também ficam ligadas. Derivando a equação em relação a $t$ (implicitamente),
obtemos uma relação entre as taxas.

\begin{estrategia}{taxas relacionadas}
\begin{enumerate}[label=\arabic*.]
\item Faça um desenho e dê nome às variáveis, indicando quais dependem do tempo.
\item Anote as taxas conhecidas e a taxa pedida, como derivadas em relação a $t$.
\item Escreva uma equação que relacione as variáveis (geometria, física etc.).
\item Derive a equação em relação a $t$, usando a regra da cadeia.
\item Só então substitua os valores do instante pedido e isole a taxa desejada.
\end{enumerate}
\end{estrategia}

\begin{nota}
Substitua os valores numéricos \emph{depois} de derivar. Se você substituir antes,
grandezas que variam viram constantes e suas derivadas somem indevidamente.
\end{nota}

\begin{exemplo}{4.1}{Enchendo um balão}
Um balão esférico é inflado a $3$ cm$^3$/s. Com que rapidez o raio cresce quando $r=8$ cm?
\solucao
$V=\frac43\pi r^3\Rightarrow\frac{dV}{dt}=4\pi r^2\frac{dr}{dt}$. Com $\frac{dV}{dt}=3$ e
$r=8$: $\frac{dr}{dt}=\frac{3}{256\pi}\approx0{,}0037$ cm/s. Quanto maior o balão, mais devagar
o raio cresce para a mesma vazão.
\end{exemplo}

\begin{exemplo}{4.2}{A escada que escorrega}
Uma escada de 10 m está apoiada numa parede. A base se afasta da parede a 1 m/s. Com que
rapidez o topo desce quando a base está a 6 m da parede?
\solucao
\figtikz{\begin{tikzpicture}[scale=0.42]
\fill[cinza!25] (-0.4,0) rectangle (0,9); \draw[thick] (0,0) -- (0,9); \draw[thick] (0,0) -- (9,0);
\draw[azul,line width=2pt] (0,8) -- (6,0);
\draw[<->,laranja] (0,-0.7) -- node[below,font=\small]{$x$} (6,-0.7);
\draw[<->,laranja] (-1,0) -- node[left,font=\small]{$y$} (-1,8);
\node[azul,font=\small] at (3.8,4.6) {10 m};
\draw[-{Stealth},thick] (6.3,0.4) -- (8,0.4); \draw[-{Stealth},thick] (0.4,8.3) -- (0.4,6.8);
\end{tikzpicture}}{A escada: $x^2+y^2=100$.}
Seja $x$ a distância da base à parede e $y$ a altura do topo: $x^2+y^2=100$. Derivando,
$2x\frac{dx}{dt}+2y\frac{dy}{dt}=0$. Quando $x=6$, $y=8$; com $\frac{dx}{dt}=1$:
$\frac{dy}{dt}=-\frac{6\cdot1}{8}=-\frac34$ m/s. O topo desce a $0{,}75$ m/s.
\end{exemplo}

\begin{exemplo}{4.3}{Enchendo um tanque cônico}
Um tanque em forma de cone invertido tem 4 m de altura e raio do topo de 2 m. Entra água a
$2$ m$^3$/min. Com que rapidez o nível sobe quando a água tem 3 m de profundidade?
\solucao
Por semelhança de triângulos, $r=\frac h2$. Então $V=\frac13\pi r^2h=\frac{\pi h^3}{12}$ e
$\frac{dV}{dt}=\frac{\pi h^2}{4}\frac{dh}{dt}$. Com $h=3$: $2=\frac{9\pi}{4}\frac{dh}{dt}$, logo
$\frac{dh}{dt}=\frac{8}{9\pi}\approx0{,}28$ m/min.
\end{exemplo}

% ---------------------------------------------------------------------
\section{Aproximações lineares e diferenciais}

Perto do ponto de tangência, o gráfico de uma função derivável é quase indistinguível de
sua reta tangente. A \textbf{aproximação linear} (ou linearização) de $f$ em $a$ é
\[
L(x)=f(a)+f'(a)(x-a),\qquad f(x)\approx L(x)\ \text{ para } x \text{ perto de } a .
\]

\begin{exemplo}{4.4}{Linearizações úteis}
(a) Para $f(x)=\sqrt x$ em $a=9$: $L(x)=3+\frac16(x-9)$, logo
$\sqrt{9{,}1}\approx3+\frac{0{,}1}6\approx3{,}0167$ (o valor exato é $3{,}01662\ldots$).

(b) Para $f(x)=\sin x$ em $a=0$: $L(x)=x$. É a aproximação $\sin x\approx x$ para ângulos
pequenos, usada, por exemplo, no estudo do pêndulo.

(c) Para $f(x)=(1+x)^n$ em $a=0$: $(1+x)^n\approx1+nx$. Assim, $\sqrt[3]{1{,}03}\approx1{,}01$.
\end{exemplo}

\textbf{Diferenciais.} Se $y=f(x)$, a \textbf{diferencial} $dy=f'(x)\,dx$ aproxima a variação
real $\Delta y=f(x+dx)-f(x)$. Na prática, as diferenciais estimam como um erro de medida se
propaga.

\begin{exemplo}{4.5}{Propagação de erro}
A aresta de um cubo mede $5$ cm, com erro possível de $\pm0{,}1$ cm. Estime o erro no volume.
\solucao
$V=x^3$, $dV=3x^2\,dx=3(25)(0{,}1)=7{,}5$ cm$^3$. O volume é $125\pm7{,}5$ cm$^3$. O erro
relativo é $\frac{dV}V=\frac{7{,}5}{125}=6\%$, o triplo do erro relativo da aresta ($2\%$),
como sugere $\frac{dV}{V}=3\frac{dx}{x}$.
\end{exemplo}

% ---------------------------------------------------------------------
\section{Máximos e mínimos}

\begin{definicao}
$f$ tem um \textbf{máximo absoluto} em $c$ se $f(c)\ge f(x)$ para todo $x$ do domínio, e um
\textbf{mínimo absoluto} se $f(c)\le f(x)$ para todo $x$. Tem um \textbf{máximo local}
(ou relativo) em $c$ se $f(c)\ge f(x)$ para todo $x$ em algum intervalo aberto contendo $c$;
analogamente para \textbf{mínimo local}. Esses valores são os \textbf{extremos} de $f$.
\end{definicao}

\figtikz{\begin{tikzpicture}
\begin{axis}[painel,width=10cm,height=5cm,xmin=-0.2,xmax=6.3,ymin=-1.5,ymax=3.8,xtick=\empty,ytick=\empty]
\addplot[domain=0.3:6,samples=200]{0.18*(x-0.3)*(x-2.4)*(x-4.6)*(x-5.6)+1.2};
\fill[laranja] (axis cs:0.3,1.2) circle (2pt);
\fill[laranja] (axis cs:6,{0.18*(5.7)*(3.6)*(1.4)*(0.4)+1.2}) circle (2pt) node[left,font=\scriptsize]{máx. absoluto};
\node[font=\scriptsize] at (axis cs:1.35,3.1) {máx. local};
\node[font=\scriptsize] at (axis cs:3.55,-1.15) {mín. absoluto (e local)};
\node[font=\scriptsize] at (axis cs:5.15,1.9) {máx. local};
\end{axis}\end{tikzpicture}}{Extremos locais e absolutos de uma função contínua em um
intervalo fechado. Os extremos absolutos podem ocorrer nas extremidades.}

\begin{teorema}[Teorema do Valor Extremo]{4.1}
Se $f$ é contínua em um intervalo fechado e limitado $[a,b]$, então $f$ tem máximo e mínimo
absolutos em $[a,b]$.
\end{teorema}

As duas hipóteses importam: $f(x)=x$ em $(0,1)$ não tem máximo nem mínimo (o intervalo é
aberto), e $f(x)=1/x$ em $[-1,1]\setminus\{0\}$ também não (falta continuidade).

\begin{teorema}[Teorema de Fermat]{4.2}
Se $f$ tem um extremo local em $c$ e $f$ é derivável em $c$, então $f'(c)=0$.
\end{teorema}

Um \textbf{ponto crítico} de $f$ é um ponto $c$ do domínio em que $f'(c)=0$ ou $f'(c)$ não
existe. Pelo teorema de Fermat, \emph{extremos locais só ocorrem em pontos críticos}. A
recíproca é falsa: $f(x)=x^3$ tem $f'(0)=0$ e nenhum extremo em $0$.

\begin{estrategia}{extremos absolutos em $[a,b]$ (método do intervalo fechado)}
\begin{enumerate}[label=\arabic*.]
\item Calcule $f$ nas extremidades $a$ e $b$.
\item Encontre os pontos críticos em $(a,b)$ e calcule $f$ neles.
\item O maior desses valores é o máximo absoluto; o menor, o mínimo absoluto.
\end{enumerate}
\end{estrategia}

\begin{exemplo}{4.6}{Extremos absolutos}
Encontre os extremos absolutos de $f(x)=x^2-3x^{2/3}$ em $[0,2]$.
\solucao
$f'(x)=2x-2x^{-1/3}=\dfrac{2(x^{4/3}-1)}{x^{1/3}}$. Pontos críticos: $f'(x)=0$ em $x=1$, e
$f'$ não existe em $x=0$ (que é também uma extremidade). Valores: $f(0)=0$, $f(1)=-2$ e
$f(2)=4-3\sqrt[3]4\approx-0{,}76$. Máximo absoluto $0$ em $x=0$; mínimo absoluto $-2$ em $x=1$.
\end{exemplo}

% ---------------------------------------------------------------------
\section{O Teorema do Valor Médio}

\begin{teorema}[Teorema de Rolle]{4.3}
Se $f$ é contínua em $[a,b]$, derivável em $(a,b)$ e $f(a)=f(b)$, então existe $c\in(a,b)$ com
$f'(c)=0$.
\end{teorema}
\emph{Ideia da prova.} Pelo Teorema do Valor Extremo, $f$ tem máximo e mínimo em $[a,b]$. Se
ambos ocorrem nas extremidades, $f$ é constante e $f'\equiv0$. Caso contrário, um deles
ocorre em um ponto interior $c$, e pelo teorema de Fermat $f'(c)=0$. $\square$

\begin{teorema}[Teorema do Valor Médio]{4.4}
Se $f$ é contínua em $[a,b]$ e derivável em $(a,b)$, então existe $c\in(a,b)$ tal que
\[
f'(c)=\frac{f(b)-f(a)}{b-a}.
\]
\end{teorema}

\figtikz{\begin{tikzpicture}
\begin{axis}[painel,width=8cm,height=5.4cm,xmin=0,xmax=5,ymin=0,ymax=4.5,xtick={0.6,4.4},xticklabels={$a$,$b$},ytick=\empty]
\addplot[domain=0.3:4.7]{0.5+0.9*x-0.08*x^2+0.6*sin(deg(1.4*x))};
\pgfmathsetmacro{\fa}{0.5+0.9*0.6-0.08*0.36+0.6*sin(deg(0.84))}
\pgfmathsetmacro{\fb}{0.5+0.9*4.4-0.08*19.36+0.6*sin(deg(6.16))}
\pgfmathsetmacro{\m}{(\fb-\fa)/3.8}
\addplot[laranja,domain=0.6:4.4]{\fa+\m*(x-0.6)};
\fill[laranja] (axis cs:0.6,\fa) circle (2pt); \fill[laranja] (axis cs:4.4,\fb) circle (2pt);
\addplot[verde,dashed,domain=1.6:4.6]{0.5+0.9*3.17-0.08*3.17^2+0.6*sin(deg(1.4*3.17))+\m*(x-3.17)};
\fill[verde] (axis cs:3.17,{0.5+0.9*3.17-0.08*3.17^2+0.6*sin(deg(1.4*3.17))}) circle (2pt) node[below,font=\scriptsize]{$c$};
\end{axis}\end{tikzpicture}}{Teorema do Valor Médio: existe um ponto $c$ em que a tangente
(verde) é paralela à secante entre as extremidades (laranja).}

Em termos físicos: se um carro percorre 200 km em 2 horas, em algum instante sua velocidade
foi exatamente 100 km/h.

\begin{exemplo}{4.7}{Verificando o TVM}
Para $f(x)=\sqrt x$ em $[0,9]$: $f$ é contínua em $[0,9]$ e derivável em $(0,9)$. A
inclinação média é $\frac{3-0}{9}=\frac13$, e $f'(c)=\frac1{2\sqrt c}=\frac13$ dá
$c=\frac94\in(0,9)$.
\end{exemplo}

O TVM é a ponte entre o sinal da derivada e o comportamento da função:

\begin{teorema}[Consequências do TVM]{4.5}
Seja $f$ derivável em um intervalo $I$.
\begin{enumerate}[label=(\roman*)]
\item Se $f'(x)=0$ para todo $x\in I$, então $f$ é constante em $I$.
\item Se $f'(x)=g'(x)$ para todo $x\in I$, então $f(x)=g(x)+C$ para alguma constante $C$.
\item Se $f'>0$ em $I$, $f$ é crescente em $I$; se $f'<0$, $f$ é decrescente.
\end{enumerate}
\end{teorema}
\emph{Prova de (iii).} Se $x_1<x_2$ em $I$, o TVM dá $c$ com
$f(x_2)-f(x_1)=f'(c)(x_2-x_1)>0$. $\square$

% ---------------------------------------------------------------------
\section{Derivadas e a forma do gráfico}

\begin{teorema}[Teste da primeira derivada]{4.6}
Seja $c$ um ponto crítico de $f$, contínua em um intervalo aberto contendo $c$.
\begin{itemize}
\item Se $f'$ muda de positiva para negativa em $c$, $f$ tem um máximo local em $c$.
\item Se $f'$ muda de negativa para positiva em $c$, $f$ tem um mínimo local em $c$.
\item Se $f'$ não muda de sinal em $c$, não há extremo local em $c$.
\end{itemize}
\end{teorema}

\textbf{Concavidade.} Se $f'$ é crescente em $I$, o gráfico de $f$ se curva para cima, e
dizemos que $f$ é \textbf{côncava para cima} em $I$. Se $f'$ é decrescente, $f$ é
\textbf{côncava para baixo}. Um \textbf{ponto de inflexão} é um ponto do gráfico em que a
concavidade muda.

\begin{teorema}[Teste da concavidade]{4.7}
Se $f''>0$ em $I$, $f$ é côncava para cima em $I$; se $f''<0$, côncava para baixo.
\end{teorema}

\begin{teorema}[Teste da segunda derivada]{4.8}
Suponha $f'(c)=0$ e $f''$ contínua perto de $c$. Se $f''(c)>0$, $f$ tem mínimo local em $c$;
se $f''(c)<0$, máximo local. Se $f''(c)=0$, o teste é inconclusivo.
\end{teorema}

\begin{exemplo}{4.8}{Análise completa de um polinômio}
Estude $f(x)=x^3-3x^2-9x+1$.
\solucao
$f'(x)=3x^2-6x-9=3(x-3)(x+1)$: crescente em $(-\infty,-1)$ e $(3,\infty)$, decrescente em
$(-1,3)$. Máximo local $f(-1)=6$; mínimo local $f(3)=-26$.
$f''(x)=6x-6$: côncava para baixo em $(-\infty,1)$ e para cima em $(1,\infty)$; inflexão em
$(1,-10)$. Pelo teste da segunda derivada, $f''(-1)=-12<0$ confirma o máximo e $f''(3)=12>0$,
o mínimo.
\figtikz{\begin{tikzpicture}
\begin{axis}[painel,width=8cm,height=5.4cm,xmin=-3,xmax=5.5,ymin=-32,ymax=14,xtick={-1,1,3},ytick={-26,-10,6}]
\addplot[domain=-3:1,samples=120]{x^3-3*x^2-9*x+1};
\addplot[laranja,domain=1:5.3,samples=120]{x^3-3*x^2-9*x+1};
\fill[black] (axis cs:-1,6) circle (2pt); \fill[black] (axis cs:3,-26) circle (2pt);
\fill[verde] (axis cs:1,-10) circle (2.2pt);
\node[font=\scriptsize,azul] at (axis cs:-2.1,-14) {côncava p/ baixo};
\node[font=\scriptsize,laranja] at (axis cs:4.2,-10) {côncava p/ cima};
\end{axis}\end{tikzpicture}}{$f(x)=x^3-3x^2-9x+1$: extremos locais (pretos) e ponto de
inflexão (verde). A cor indica a concavidade.}
\end{exemplo}

% ---------------------------------------------------------------------
\section{Limites no infinito e assíntotas}

Escrevemos $\lim_{x\to\infty}f(x)=L$ se $f(x)$ fica arbitrariamente perto de $L$ para $x$
suficientemente grande; analogamente para $x\to-\infty$. Nesse caso, a reta $y=L$ é uma
\textbf{assíntota horizontal}. As propriedades dos limites continuam valendo, e o limite
básico é $\lim_{x\to\pm\infty}\frac1{x^n}=0$ para $n>0$.

\textbf{Funções racionais.} Dividimos numerador e denominador pela maior potência de $x$ do
denominador. Por exemplo,
\[
\lim_{x\to\infty}\frac{3x^2+5}{2x^2-x}=\lim_{x\to\infty}\frac{3+5/x^2}{2-1/x}=\frac32 .
\]
Em geral, para $\frac{p(x)}{q(x)}$: se $\operatorname{grau}p<\operatorname{grau}q$, o limite
é $0$; se os graus são iguais, é a razão dos coeficientes dominantes; se
$\operatorname{grau}p>\operatorname{grau}q$, o limite é $\pm\infty$. Quando
$\operatorname{grau}p=\operatorname{grau}q+1$, a divisão de polinômios dá uma
\textbf{assíntota oblíqua}: $\frac{p(x)}{q(x)}=mx+b+\frac{r(x)}{q(x)}$, com
$\frac{r}{q}\to0$.

\begin{nota}
Cuidado com raízes quando $x\to-\infty$: $\sqrt{x^2}=|x|=-x$. Assim,
$\lim_{x\to\infty}\frac{\sqrt{x^2+1}}{x}=1$, mas $\lim_{x\to-\infty}\frac{\sqrt{x^2+1}}{x}=-1$.
\end{nota}

Comportamentos transcendentes a lembrar: $e^x\to\infty$ e $e^{-x}\to0$ quando $x\to\infty$;
$\ln x\to\infty$ (devagar); $\tan^{-1}x\to\pm\frac\pi2$ quando $x\to\pm\infty$.

\begin{estrategia}{esboço do gráfico de uma função}
\begin{enumerate}[label=\arabic*.]
\item Domínio, interceptos e simetria.
\item Assíntotas verticais (limites infinitos) e horizontais/oblíquas (limites no infinito).
\item $f'$: intervalos de crescimento e extremos locais.
\item $f''$: concavidade e pontos de inflexão.
\item Junte tudo num esboço.
\end{enumerate}
\end{estrategia}

\begin{exemplo}{4.9}{Esboço com assíntota oblíqua}
Esboce $f(x)=\dfrac{x^2}{x-1}$.
\solucao
\emph{Domínio:} $x\neq1$; intercepto $(0,0)$. \emph{Assíntotas:} vertical $x=1$
($-\infty$ pela esquerda, $+\infty$ pela direita); como $\frac{x^2}{x-1}=x+1+\frac1{x-1}$, a
reta $y=x+1$ é assíntota oblíqua. \emph{Primeira derivada:}
$f'(x)=\frac{x(x-2)}{(x-1)^2}$: máximo local $f(0)=0$, mínimo local $f(2)=4$.
\emph{Segunda derivada:} $f''(x)=\frac2{(x-1)^3}$: côncava para baixo em $x<1$, para cima em
$x>1$, sem inflexão (que exigiria $x=1$, fora do domínio).
\figtikz{\begin{tikzpicture}
\begin{axis}[painel,width=7.5cm,height=6cm,xmin=-4,xmax=6,ymin=-8,ymax=11,xtick={1,2},ytick={4}]
\addplot[domain=-4:0.85,samples=200]{x^2/(x-1)}; \addplot[domain=1.12:6,samples=200]{x^2/(x-1)};
\addplot[laranja,dashed,domain=-4:6]{x+1}; \draw[laranja,dashed] (axis cs:1,-8) -- (axis cs:1,11);
\fill[black] (axis cs:0,0) circle (2pt); \fill[black] (axis cs:2,4) circle (2pt);
\end{axis}\end{tikzpicture}}{$f(x)=x^2/(x-1)$, com assíntotas $x=1$ e $y=x+1$.}
\end{exemplo}

% ---------------------------------------------------------------------
\section{Otimização}

Muitos problemas práticos pedem o maior ou o menor valor de uma grandeza: área máxima, custo
mínimo, tempo mínimo. A estratégia combina modelagem com o método do intervalo fechado ou com
o teste da primeira derivada.

\begin{estrategia}{problemas de otimização}
\begin{enumerate}[label=\arabic*.]
\item Faça um desenho e nomeie as variáveis.
\item Escreva a grandeza a otimizar em função das variáveis.
\item Use as restrições do problema para reduzir a uma função de \emph{uma} variável, e
  determine o intervalo de valores possíveis.
\item Encontre os pontos críticos e compare os valores (extremidades incluídas), ou use o
  teste da primeira ou da segunda derivada.
\item Responda à pergunta, com unidades.
\end{enumerate}
\end{estrategia}

\begin{exemplo}{4.10}{A caixa de volume máximo}
De uma folha de $24\times36$ cm cortam-se quadrados de lado $x$ nos cantos e dobram-se as
abas para formar uma caixa sem tampa. Qual $x$ maximiza o volume?
\solucao
$V(x)=x(36-2x)(24-2x)=4x^3-120x^2+864x$, com $0\le x\le12$. Então
$V'(x)=12x^2-240x+864=12(x^2-20x+72)=0$ dá $x=10\pm\sqrt{28}$; só $x=10-2\sqrt7\approx4{,}71$
está em $[0,12]$. Como $V(0)=V(12)=0$, o máximo é $V(4{,}71)\approx1825$ cm$^3$.
\end{exemplo}

\begin{exemplo}{4.11}{A lata mais econômica}
Uma lata cilíndrica fechada deve ter volume $V$. Que proporção entre altura e raio minimiza a
área de material?
\solucao
$V=\pi r^2h\Rightarrow h=\frac V{\pi r^2}$. A área é
$A(r)=2\pi r^2+2\pi rh=2\pi r^2+\frac{2V}r$, para $r>0$. $A'(r)=4\pi r-\frac{2V}{r^2}=0$ dá
$r^3=\frac V{2\pi}$. Como $A''(r)=4\pi+\frac{4V}{r^3}>0$, é mínimo. Nesse caso
$h=\frac{V}{\pi r^2}=\frac{2\pi r^3}{\pi r^2}=2r$: a lata ótima tem \emph{altura igual ao
diâmetro}.
\end{exemplo}

% ---------------------------------------------------------------------
\section{A regra de L'Hôpital}

\begin{teorema}[Regra de L'Hôpital]{4.9}
Sejam $f$ e $g$ deriváveis perto de $a$ (exceto possivelmente em $a$), com $g'(x)\neq0$
perto de $a$. Se $\lim_{x\to a}f(x)=\lim_{x\to a}g(x)=0$, ou se ambos são $\pm\infty$, então
\[
\lim_{x\to a}\frac{f(x)}{g(x)}=\lim_{x\to a}\frac{f'(x)}{g'(x)},
\]
desde que o limite da direita exista (ou seja $\pm\infty$). O mesmo vale para $x\to\pm\infty$ e
para limites laterais.
\end{teorema}

\begin{nota}
Deriva-se numerador e denominador \emph{separadamente}; não se usa a regra do quociente. E
a regra só vale para as formas $\frac00$ e $\frac\infty\infty$: por exemplo,
$\lim_{x\to0}\frac{x+1}{x+2}=\frac12$, enquanto ``aplicar L'Hôpital'' daria, erradamente, 1.
\end{nota}

\begin{exemplo}{4.12}{As formas $\frac00$ e $\frac\infty\infty$}
(a) $\displaystyle\lim_{x\to0}\frac{1-\cos x}{x^2}=\lim_{x\to0}\frac{\sin x}{2x}=\frac12$.\qquad
(b) $\displaystyle\lim_{x\to\infty}\frac{e^x}{x^2}=\lim_{x\to\infty}\frac{e^x}{2x}
=\lim_{x\to\infty}\frac{e^x}{2}=\infty$.
\end{exemplo}

\textbf{Outras formas indeterminadas.} As formas $0\cdot\infty$ e $\infty-\infty$ são
reescritas como quocientes; as formas $1^\infty$, $0^0$ e $\infty^0$ são tratadas com
logaritmo: se $y=f(x)^{g(x)}$, calculamos $\lim\ln y=\lim g(x)\ln f(x)$ e depois
exponenciamos.

\begin{exemplo}{4.13}{Formas indeterminadas diversas}
(a) $0\cdot\infty$: $\displaystyle\lim_{x\to0^+}x\ln x=\lim_{x\to0^+}\frac{\ln x}{1/x}
=\lim_{x\to0^+}\frac{1/x}{-1/x^2}=\lim_{x\to0^+}(-x)=0$.

(b) $\infty-\infty$: $\displaystyle\lim_{x\to0^+}\Big(\frac1x-\frac1{\sin x}\Big)
=\lim_{x\to0^+}\frac{\sin x-x}{x\sin x}=\lim_{x\to0^+}\frac{\cos x-1}{\sin x+x\cos x}
=\lim_{x\to0^+}\frac{-\sin x}{2\cos x-x\sin x}=0$.

(c) $1^\infty$: para $y=\big(1+\frac1x\big)^x$, $\ln y=x\ln\big(1+\frac1x\big)
=\frac{\ln(1+1/x)}{1/x}\to1$ (por L'Hôpital). Logo $\lim_{x\to\infty}\big(1+\frac1x\big)^x=e$.

(d) $0^0$: para $y=x^x$, $\ln y=x\ln x\to0$ por (a). Logo $\lim_{x\to0^+}x^x=e^0=1$.
\end{exemplo}

\textbf{Taxas de crescimento.} Aplicando L'Hôpital repetidamente: para todo $p>0$,
$\frac{\ln x}{x^p}\to0$ e $\frac{x^p}{e^x}\to0$ quando $x\to\infty$. Ou seja, no infinito,
\emph{o logaritmo cresce mais devagar que qualquer potência, e toda potência cresce mais
devagar que a exponencial}.

% ---------------------------------------------------------------------
\section{O método de Newton}

Para aproximar uma raiz de $f(x)=0$, partimos de um palpite $x_0$ e trocamos a curva pela
tangente em $(x_0,f(x_0))$. O ponto em que a tangente corta o eixo $x$ é o novo palpite:
\[
x_{n+1}=x_n-\frac{f(x_n)}{f'(x_n)} .
\]

\figtikz{\begin{tikzpicture}
\begin{axis}[painel,width=8cm,height=5.2cm,xmin=0.8,xmax=3.2,ymin=-1.5,ymax=7.5,xtick={3,1.8333,1.4621},xticklabels={$x_0$,$x_1$,$x_2$},ytick=\empty,
  x tick label style={font=\scriptsize}]
\addplot[domain=0.8:3.2]{x^2-2};
\addplot[laranja,domain=1.7:3.2]{7+6*(x-3)};
\addplot[laranja!70,domain=1.35:2.3]{1.3611+3.6667*(x-1.8333)};
\draw[cinza,dotted] (axis cs:3,0) -- (axis cs:3,7); \draw[cinza,dotted] (axis cs:1.8333,0) -- (axis cs:1.8333,1.3611);
\fill[verde] (axis cs:1.4142,0) circle (2pt) node[below left,font=\scriptsize]{$\sqrt2$};
\end{axis}\end{tikzpicture}}{Método de Newton para $f(x)=x^2-2$: cada tangente leva a um
palpite mais próximo da raiz $\sqrt2$.}

\begin{exemplo}{4.14}{Calculando $\sqrt2$}
Para $f(x)=x^2-2$, a iteração é $x_{n+1}=x_n-\frac{x_n^2-2}{2x_n}=\frac12\Big(x_n+\frac2{x_n}\Big)$.
Com $x_0=1$: $x_1=1{,}5$; $x_2\approx1{,}416667$; $x_3\approx1{,}414216$; $x_4\approx1{,}4142136$.
O número de casas corretas praticamente dobra a cada passo.
\end{exemplo}

\textbf{Quando falha.} Se $f'(x_n)=0$, a tangente é horizontal e o método para. Ele também
pode oscilar: para $f(x)=x^3-2x+2$ com $x_0=0$, obtemos $x_1=1$, $x_2=0$, $x_3=1,\dots$, sem
convergir. Por isso convém escolher $x_0$ perto da raiz, por exemplo com o auxílio de um
gráfico ou do Teorema do Valor Intermediário.

% ---------------------------------------------------------------------
\section{Antiderivadas}

\begin{definicao}
$F$ é uma \textbf{antiderivada} (ou primitiva) de $f$ em um intervalo $I$ se $F'(x)=f(x)$ para
todo $x\in I$.
\end{definicao}

Pela consequência (ii) do TVM, se $F$ é uma antiderivada de $f$, então \emph{todas} as
antiderivadas são da forma $F(x)+C$. Escrevemos a \textbf{integral indefinida}
\[
\int f(x)\,dx=F(x)+C .
\]

\figtikz{\begin{tikzpicture}
\begin{axis}[painel,width=6.5cm,height=5cm,xmin=-2.3,xmax=2.3,ymin=-2.5,ymax=5,xtick=\empty,ytick=\empty]
\foreach \c/\o in {-2/40,-1/55,0/100,1/55,2/40}{\edef\tmp{\noexpand\addplot[azul!\o!white,domain=-2.3:2.3]{x^2+\c};}\tmp}
\end{axis}\end{tikzpicture}}{A família $x^2+C$: todas as antiderivadas de $2x$ diferem por
uma translação vertical.}

Cada regra de derivação, lida ao contrário, dá uma regra de antiderivação:
\[
\int x^n\,dx=\frac{x^{n+1}}{n+1}+C\ (n\neq-1),\qquad
\int\frac1x\,dx=\ln|x|+C,\qquad \int e^x\,dx=e^x+C,
\]
\[
\int\cos x\,dx=\sin x+C,\quad \int\sin x\,dx=-\cos x+C,\quad \int\sec^2x\,dx=\tan x+C,\quad
\int\frac{dx}{1+x^2}=\tan^{-1}x+C .
\]
Além disso, $\int\big(kf\pm g\big)\,dx=k\int f\,dx\pm\int g\,dx$. Uma antiderivada sempre pode
ser conferida derivando.

\begin{exemplo}{4.15}{Calculando integrais indefinidas}
$\displaystyle\int\Big(5x^3-\frac{7}{x^2}+4\sqrt x\Big)dx
=\frac{5x^4}4+\frac7x+\frac{8}{3}x^{3/2}+C$
\quad e\quad
$\displaystyle\int\frac{x^2+1}{x}\,dx=\int\Big(x+\frac1x\Big)dx=\frac{x^2}2+\ln|x|+C$.
\end{exemplo}

\textbf{Problemas de valor inicial.} Uma condição como $F(a)=b$ fixa a constante $C$.

\begin{exemplo}{4.16}{Um carro freando}
Um carro a 30 m/s freia com desaceleração constante de 6 m/s$^2$. Que distância percorre até
parar?
\solucao
$a(t)=-6\Rightarrow v(t)=-6t+C_1$, e $v(0)=30$ dá $v(t)=30-6t$. Então
$s(t)=30t-3t^2+C_2$, com $s(0)=0$. O carro para quando $v=0$, em $t=5$ s, tendo percorrido
$s(5)=150-75=75$ m.
\end{exemplo}

% ---------------------------------------------------------------------
\section*{Resumo do capítulo}
\addcontentsline{toc}{section}{Resumo do capítulo}
\begin{itemize}
\item Taxas relacionadas: derive a equação que liga as variáveis em relação a $t$; substitua
  os valores só depois.
\item Linearização $L(x)=f(a)+f'(a)(x-a)$; diferencial $dy=f'(x)\,dx$ estima variações e
  erros.
\item Extremos locais ocorrem em pontos críticos ($f'=0$ ou $\nexists$). Em $[a,b]$ com $f$
  contínua, compare $f$ nos pontos críticos e nas extremidades.
\item TVM: $f'(c)=\frac{f(b)-f(a)}{b-a}$. Consequências: $f'=0\Rightarrow f$ constante; sinal
  de $f'$ dá crescimento.
\item Testes da 1ª e da 2ª derivada; $f''$ dá a concavidade; inflexão onde ela muda.
\item Assíntotas horizontais e oblíquas via limites no infinito; roteiro de esboço.
\item Otimização: modele, reduza a uma variável, encontre e compare pontos críticos.
\item L'Hôpital para $\frac00$ e $\frac\infty\infty$; as demais formas se reduzem a essas.
  $\ln x\ll x^p\ll e^x$.
\item Newton: $x_{n+1}=x_n-f(x_n)/f'(x_n)$.
\item Antiderivadas: $\int f\,dx=F+C$; problemas de valor inicial fixam $C$.
\end{itemize}

% Gerado por ferramentas/gerar.py; não edite à mão.
\section*{Exercícios}
\addcontentsline{toc}{section}{Exercícios}
Os pontos antes de cada número indicam a dificuldade: \nivel{1} fácil, \nivel{2} médio, \nivel{3} difícil e \nivel{4} desafio. Os exercícios marcados com \estrela{} têm solução completa, com animação do gráfico, no repositório do livro: \url{https://github.com/AleffAzevedo/calculo-1-pt/blob/main/exercicios/cap4/README.md}. As respostas estão no fim do capítulo.

\begin{itemize}[leftmargin=5.2em,labelsep=0.6em,itemsep=4pt]
\item[\exnum{1}{4.1}] Encontre a linearização de $f(x)=\sqrt[3]x$ em $a=8$ e use-a para estimar $\sqrt[3]{8{,}2}$.
\item[\exnum{1}{4.2}] Para $y=x^2-3x$, compare $dy$ e $\Delta y$ quando $x=2$ e $dx=0{,}1$.
\item[\exnum{1}{4.3}] Encontre os pontos críticos de $f(x)=x^3-12x+1$.
\item[\exnum{1}{4.4}] Encontre os pontos críticos de $f(x)=x^{2/3}(x-5)$ (inclua os pontos em que $f'$ não existe).
\item[\exnum{1}{4.5}] Encontre os extremos absolutos de $f(x)=x^2-4x+3$ em $[0,3]$.
\item[\exnum{1}{4.6}] Encontre os extremos absolutos de $f(x)=\sin x+\cos x$ em $[0,\pi]$.
\item[\exnum{1}{4.7}] Verifique as hipóteses do Teorema de Rolle para $f(x)=x^2-4x$ em $[0,4]$ e encontre $c$.
\item[\exnum{1}{4.8}] Encontre o $c$ do TVM para $f(x)=x^2$ em $[1,3]$ e interprete.
\item[\exnum{1}{4.9}] Determine os intervalos de crescimento e de decrescimento de $f(x)=x^3-3x^2$.
\item[\exnum{1}{4.10}] Estude a concavidade de $f(x)=x^4-4x^3$ e encontre os pontos de inflexão.
\item[\exnum{1}{4.11}] Use o teste da segunda derivada para classificar os pontos críticos de $f(x)=2x^3-3x^2-12x$.
\item[\exnum{1}{4.12}] Calcule $\displaystyle\lim_{x\to\infty}\frac{3x^2-x}{5x^2+2}$.
\item[\exnum{1}{4.13}] Calcule $\displaystyle\lim_{x\to-\infty}\frac{2x+1}{\sqrt{x^2+3}}$.
\item[\exnum{1}{4.14}] Encontre as assíntotas verticais e horizontais de $f(x)=\dfrac{2x^2}{x^2-9}$.
\item[\exnum{1}{4.15}] Encontre a assíntota oblíqua de $f(x)=\dfrac{x^2+1}{x+1}$.
\item[\exnum{1}{4.16}] Use L'Hôpital para calcular $\displaystyle\lim_{x\to1}\frac{\ln x}{x-1}$ e $\displaystyle\lim_{x\to0}\frac{e^x-1}{x}$.
\item[\exnum{1}{4.17}] Calcule $\displaystyle\lim_{x\to\infty}x^2e^{-x}$.
\item[\exnum{1}{4.18}] Faça duas iterações do método de Newton para $f(x)=x^2-5$ com $x_0=2$.
\item[\exnum{1}{4.19}] Calcule $\int(3x^2-4x+5)\,dx$ e $\int(\sin x+e^x)\,dx$.
\item[\exnum{1}{4.20}] Resolva o problema de valor inicial $f'(x)=6x^2-2$, $f(1)=3$.
\item[\exnum{2}{4.21}] \estrela\ \textbf{A escada escorregando.} Uma escada de 5 m está apoiada numa parede vertical. A base se afasta da parede a
$0{,}5$ m/s. (a) Com que rapidez o topo desce quando a base está a 3 m da parede?
(b) Em que posição as duas extremidades se movem com a mesma rapidez?
\item[\exnum{2}{4.22}] \estrela\ \textbf{A sombra no poste.} Uma pessoa de $1{,}8$ m se afasta de um poste de luz de 6 m de altura a $1{,}5$ m/s.
(a) Com que rapidez sua sombra aumenta? (b) Com que rapidez a ponta da sombra se move?
\item[\exnum{2}{4.23}] \estrela\ \textbf{Avião e radar.} Um avião voa horizontalmente a 4 km de altitude, a 600 km/h, numa trajetória que passa
diretamente acima de uma estação de radar. Com que rapidez a distância avião–radar aumenta
quando essa distância é 5 km (e o avião já passou da estação)?
\item[\exnum{2}{4.24}] Uma mancha de óleo circular se espalha com o raio crescendo a 2 m/s. Com que rapidez a área cresce quando $r=10$ m?
\item[\exnum{2}{4.25}] A aresta de um cubo cresce a $0{,}5$ cm/s. Com que rapidez o volume cresce quando a aresta mede 4 cm?
\item[\exnum{2}{4.26}] O raio de uma esfera mede $10$ cm com erro de $\pm0{,}2$ cm. Estime o erro absoluto e o percentual no volume.
\item[\exnum{2}{4.27}] Use uma aproximação linear para estimar $\sin31^\circ$.
\item[\exnum{2}{4.28}] \estrela\ \textbf{Método do intervalo fechado.} Encontre os valores máximo e mínimo absolutos de $f(x)=x\,e^{-x}$ em $[0,4]$.
\item[\exnum{2}{4.29}] \estrela\ \textbf{Teorema do Valor Médio.} Verifique as hipóteses do TVM para $f(x)=x^3-x$ em $[0,2]$ e encontre todos os $c$ garantidos
pelo teorema.
\item[\exnum{2}{4.30}] Mostre que $x^3+x-1=0$ tem exatamente uma raiz real.
\item[\exnum{2}{4.31}] Mostre que $\tan^{-1}x+\tan^{-1}\tfrac1x=\tfrac\pi2$ para $x>0$.
\item[\exnum{2}{4.32}] Encontre os extremos locais de $f(x)=x^2e^{-x}$.
\item[\exnum{2}{4.33}] Encontre os extremos locais de $f(x)=x-2\sin x$ em $(0,2\pi)$.
\item[\exnum{2}{4.34}] Calcule $\displaystyle\lim_{x\to\infty}\big(\sqrt{x^2+x}-x\big)$.
\item[\exnum{2}{4.35}] Calcule $\displaystyle\lim_{x\to0}\frac{x-\sin x}{x^3}$.
\item[\exnum{2}{4.36}] Calcule $\displaystyle\lim_{x\to\infty}x^{1/x}$.
\item[\exnum{2}{4.37}] Calcule $\displaystyle\lim_{x\to\pi/2^-}(\sec x-\tan x)$.
\item[\exnum{2}{4.38}] \estrela\ \textbf{O limite $(1+a/x)^x$.} Mostre, com L'Hôpital, que $\displaystyle\lim_{x\to\infty}\Big(1+\frac ax\Big)^x=e^a$ para
qualquer constante $a$. Use o resultado para calcular o montante de R\$\,1000 a 6\% ao ano
com capitalização contínua por 10 anos.
\item[\exnum{2}{4.39}] Calcule $\displaystyle\int\frac{x^3-2x}{x^2}\,dx$.
\item[\exnum{2}{4.40}] Calcule $\displaystyle\int\Big(\frac2{1+x^2}-\frac3{\sqrt{1-x^2}}\Big)dx$.
\item[\exnum{3}{4.41}] \estrela\ \textbf{Esboço completo de um gráfico.} Faça o estudo completo e esboce $f(x)=\dfrac{x}{x^2+1}$: domínio, simetria, assíntotas,
crescimento, extremos, concavidade e inflexões.
\item[\exnum{3}{4.42}] \estrela\ \textbf{Cercado junto ao rio.} Um fazendeiro tem 600 m de cerca para construir um cercado retangular junto a um rio reto
(sem cerca do lado do rio), dividido em dois currais por uma cerca perpendicular ao rio. Que
dimensões maximizam a área total?
\item[\exnum{3}{4.43}] \estrela\ \textbf{Ponto mais próximo.} Encontre os pontos da parábola $y=x^2$ mais próximos do ponto $(0,3)$ e a distância mínima.
\item[\exnum{3}{4.44}] \estrela\ \textbf{O cabo mais barato.} Uma usina fica numa margem de um rio de 1 km de largura. Uma fábrica fica na outra
margem, 3 km rio abaixo. O cabo custa R\$\,5 mil/km sob a água e R\$\,4 mil/km em terra.
Qual é o trajeto mais barato (atravessar em diagonal até um ponto e seguir pela margem)?
\item[\exnum{3}{4.45}] \estrela\ \textbf{Método de Newton para $\cos x=x$.} Use o método de Newton com $x_0=1$ para aproximar a solução de $\cos x=x$ com 9 casas
decimais. Quantas iterações são necessárias?
\item[\exnum{3}{4.46}] \estrela\ \textbf{Lançamento vertical.} Uma bola é lançada para cima a partir de 2 m de altura com velocidade de 15 m/s
($g=9{,}8$ m/s²). Partindo de $a(t)=-9{,}8$, encontre $v(t)$ e $s(t)$, a altura máxima e o
instante em que a bola atinge o chão.
\item[\exnum{3}{4.47}] Uma calha de 10 m tem seção em triângulo isósceles invertido, com 1 m de largura no topo e $0{,}5$ m de altura. Entra água a $0{,}4$ m³/min. Com que rapidez o nível sobe quando $h=0{,}2$ m?
\item[\exnum{3}{4.48}] Dois carros partem do mesmo ponto: um para o norte a 60 km/h e outro para o leste a 80 km/h. Com que rapidez a distância entre eles aumenta?
\item[\exnum{3}{4.49}] Um balão sobe verticalmente a 2 m/s e é observado de um ponto no solo a 50 m da base. Com que rapidez o ângulo de elevação aumenta quando o balão está a 50 m de altura?
\item[\exnum{3}{4.50}] Encontre o retângulo de maior área inscrito num semicírculo de raio $R$ (com um lado sobre o diâmetro).
\item[\exnum{3}{4.51}] Uma caixa sem tampa de base quadrada deve ter volume $32\,000$ cm³. Que dimensões minimizam o material?
\item[\exnum{3}{4.52}] Encontre o retângulo de maior área com base no eixo $x$ e vértices superiores na parábola $y=12-x^2$.
\item[\exnum{3}{4.53}] A demanda de um produto é $q=1000-20p$ ($p$ em reais). Que preço maximiza a receita $R=pq$?
\item[\exnum{3}{4.54}] Faça o estudo completo de $f(x)=x^{2/3}(x-5)$ (há uma cúspide) e esboce o gráfico.
\item[\exnum{3}{4.55}] Faça o estudo completo de $f(x)=e^{-x^2}$ e esboce o gráfico.
\item[\exnum{3}{4.56}] Calcule $\displaystyle\lim_{x\to0}\frac{e^x-1-x}{x^2}$ e $\displaystyle\lim_{x\to\infty}x\Big(\frac\pi2-\tan^{-1}x\Big)$.
\item[\exnum{3}{4.57}] Use o TVM para provar que $|\sin a-\sin b|\le|a-b|$ para todos $a,b$.
\item[\exnum{3}{4.58}] Prove que $e^x\ge1+x$ para todo $x$.
\item[\exnum{3}{4.59}] Resolva $f''(x)=6x+2$, $f'(0)=1$, $f(0)=-1$.
\item[\exnum{3}{4.60}] Prove que a equação $\cos x=x$ tem exatamente uma solução.
\item[\exnum{4}{4.61}] Prove que a única função derivável com $f'=f$ e $f(0)=1$ é $f(x)=e^x$. (Dica: derive $g(x)=f(x)e^{-x}$.)
\item[\exnum{4}{4.62}] A luz vai de $A$ (no ar, velocidade $v_1$) a $B$ (na água, velocidade $v_2$) pelo caminho de tempo mínimo. Deduza a lei de Snell $\dfrac{\sin\theta_1}{v_1}=\dfrac{\sin\theta_2}{v_2}$.
\item[\exnum{4}{4.63}] Para $a>0$, mostre que o método de Newton para $x^2=a$ dá $x_{n+1}=\tfrac12\big(x_n+\tfrac a{x_n}\big)$ e que $x_n\ge\sqrt a$ para todo $n\ge1$, qualquer que seja $x_0>0$.
\item[\exnum{4}{4.64}] Calcule $\displaystyle\lim_{x\to0}\Big(\frac1{x^2}-\frac1{\sin^2x}\Big)$.
\item[\exnum{4}{4.65}] Qual é o comprimento da maior barra que pode ser carregada horizontalmente pela esquina entre dois corredores de larguras 8 m e 27 m?
\end{itemize}
\subsection*{Respostas}
{\small\begin{multicols}{2}\raggedright
\begin{itemize}[leftmargin=2.7em,labelsep=0.5em,itemsep=1pt]
\item[\textbf{4.1}] $L(x)=2+\tfrac{x-8}{12}$; $\approx2{,}0167$.
\item[\textbf{4.2}] $dy=0{,}1$; $\Delta y=0{,}11$.
\item[\textbf{4.3}] $x=\pm2$.
\item[\textbf{4.4}] $x=2$ e $x=0$.
\item[\textbf{4.5}] Máximo $3$ em $x=0$; mínimo $-1$ em $x=2$.
\item[\textbf{4.6}] Máximo $\sqrt2$ em $x=\tfrac\pi4$; mínimo $-1$ em $x=\pi$.
\item[\textbf{4.7}] $c=2$.
\item[\textbf{4.8}] $c=2$: tangente paralela à secante.
\item[\textbf{4.9}] Cresce em $(-\infty,0)\cup(2,\infty)$; decresce em $(0,2)$.
\item[\textbf{4.10}] Para cima em $(-\infty,0)\cup(2,\infty)$, para baixo em $(0,2)$; inflexões $(0,0)$ e $(2,-16)$.
\item[\textbf{4.11}] Máx. local $f(-1)=7$; mín. local $f(2)=-20$.
\item[\textbf{4.12}] $\tfrac35$.
\item[\textbf{4.13}] $-2$.
\item[\textbf{4.14}] Verticais $x=\pm3$; horizontal $y=2$.
\item[\textbf{4.15}] $y=x-1$.
\item[\textbf{4.16}] $1$ e $1$.
\item[\textbf{4.17}] $0$.
\item[\textbf{4.18}] $x_1=2{,}25$; $x_2\approx2{,}236111$.
\item[\textbf{4.19}] $x^3-2x^2+5x+C$; $-\cos x+e^x+C$.
\item[\textbf{4.20}] $f(x)=2x^3-2x+3$.
\item[\textbf{4.21}] (a) $0{,}375$ m/s; (b) quando $x=y=\tfrac{5}{\sqrt2}\approx3{,}54$ m.
\item[\textbf{4.22}] (a) $\tfrac{9}{14}\approx0{,}64$ m/s; (b) $\tfrac{15}{7}\approx2{,}14$ m/s.
\item[\textbf{4.23}] 360 km/h.
\item[\textbf{4.24}] $40\pi\approx125{,}7$ m²/s.
\item[\textbf{4.25}] $24$ cm³/s.
\item[\textbf{4.26}] $dV=80\pi\approx251$ cm³; $6\%$.
\item[\textbf{4.27}] $\tfrac12+\tfrac{\sqrt3}2\cdot\tfrac\pi{180}\approx0{,}5151$.
\item[\textbf{4.28}] Máximo $\tfrac1e\approx0{,}368$ em $x=1$; mínimo $0$ em $x=0$.
\item[\textbf{4.29}] $c=\tfrac{2}{\sqrt3}\approx1{,}155$.
\item[\textbf{4.30}] Existe pelo TVI; é única porque $f'=3x^2+1>0$ (pelo Rolle, duas raízes exigiriam $f'=0$).
\item[\textbf{4.31}] A derivada da soma é $0$, então ela é constante; em $x=1$ vale $\tfrac\pi2$.
\item[\textbf{4.32}] Mín. local $f(0)=0$; máx. local $f(2)=\tfrac4{e^2}\approx0{,}541$.
\item[\textbf{4.33}] Mín. $f(\tfrac\pi3)=\tfrac\pi3-\sqrt3$; máx. $f(\tfrac{5\pi}3)=\tfrac{5\pi}3+\sqrt3$.
\item[\textbf{4.34}] $\tfrac12$.
\item[\textbf{4.35}] $\tfrac16$.
\item[\textbf{4.36}] $1$.
\item[\textbf{4.37}] $0$.
\item[\textbf{4.38}] $e^a$; $1000e^{0{,}6}\approx$ R\$\,1822{,}12.
\item[\textbf{4.39}] $\tfrac{x^2}2-2\ln|x|+C$.
\item[\textbf{4.40}] $2\tan^{-1}x-3\sin^{-1}x+C$.
\item[\textbf{4.41}] Ímpar; assíntota $y=0$; máx. $(1,\tfrac12)$, mín. $(-1,-\tfrac12)$; inflexões em $x=0$ e $x=\pm\sqrt3$.
\item[\textbf{4.42}] 100 m (perpendicular ao rio) por 300 m (paralelo); área 30\,000 m².
\item[\textbf{4.43}] $\big(\pm\sqrt{5/2},\,\tfrac52\big)$; distância $\tfrac{\sqrt{11}}{2}\approx1{,}658$.
\item[\textbf{4.44}] Atravessar até o ponto a $\tfrac43$ km rio abaixo; custo mínimo R\$\,15 mil.
\item[\textbf{4.45}] $x\approx0{,}739085133$; 3 iterações (a 4ª só confirma).
\item[\textbf{4.46}] $s(t)=-4{,}9t^2+15t+2$; altura máx. $\approx13{,}48$ m em $t\approx1{,}53$ s; chão em $t\approx3{,}19$ s.
\item[\textbf{4.47}] $0{,}1$ m/min.
\item[\textbf{4.48}] $100$ km/h (constante).
\item[\textbf{4.49}] $0{,}02$ rad/s.
\item[\textbf{4.50}] Lados $\sqrt2R$ e $\tfrac R{\sqrt2}$; área $R^2$.
\item[\textbf{4.51}] Base $40\times40$ cm, altura $20$ cm; área $4800$ cm².
\item[\textbf{4.52}] Base 4, altura 8; área 32.
\item[\textbf{4.53}] $p=25$; receita R\$\,12\,500.
\item[\textbf{4.54}] Máx. local $(0,0)$ (cúspide); mín. local $(2,-3\sqrt[3]4)\approx(2;-4{,}76)$; inflexão $(-1,-6)$.
\item[\textbf{4.55}] Par; máx. $(0,1)$; inflexões em $x=\pm\tfrac1{\sqrt2}$; assíntota $y=0$.
\item[\textbf{4.56}] $\tfrac12$ e $1$.
\item[\textbf{4.57}] $\sin a-\sin b=\cos c\,(a-b)$ e $|\cos c|\le1$.
\item[\textbf{4.58}] $g(x)=e^x-1-x$ tem mínimo absoluto $g(0)=0$.
\item[\textbf{4.59}] $f(x)=x^3+x^2+x-1$.
\item[\textbf{4.60}] Existência pelo TVI; unicidade porque $g(x)=\cos x-x$ tem $g'=-\sin x-1\le0$, zerando só em pontos isolados.
\item[\textbf{4.61}] $g'=0$, logo $g$ é constante $=g(0)=1$.
\item[\textbf{4.62}] Escreva o tempo em função do ponto de cruzamento e anule a derivada.
\item[\textbf{4.63}] Desigualdade entre as médias aritmética e geométrica.
\item[\textbf{4.64}] $-\tfrac13$.
\item[\textbf{4.65}] $13\sqrt{13}\approx46{,}9$ m (minimizando $\tfrac8{\sin\theta}+\tfrac{27}{\cos\theta}$).
\end{itemize}\end{multicols}}

